\section{RQ4 \emph{(pre-registered as exploratory)}: What the Endpoint Removes Before a
Score Exists}
\label{sec:census}

\textbf{Eight of the 21 model-platform pairs we attempted never produced a score; five of the 18 on the
two platforms the design measures.} The third platform is Azure, probed while assembling the sample and
then left out of the arms. This is the result that does not depend on a $p$-value --- an endpoint either
returned the message we quote or it did not --- and it is reported at a scale a reader can size.

\textbf{Three quantities are counted here, each in its own unit and with its own denominator.}
\emph{(A)} Pairs removed before measurement: 8 of 21 attempted, spanning
11 distinct models, traceable to 7 distinct causes --- six of them are the mechanisms M1--M6 below and the
seventh is a reasoning-channel case the deposit carries. \emph{(B)} Losses inside a cell that had already
started: M7 cost 453 rows of one cell, discarded and re-collected; M8 rejected 11 trajectories across the
two retained collections --- four in the original collection of 5{,}805 valid measurements, seven in the
re-collection of 5{,}809. \emph{(C)} Distinct mechanisms
documented: 8, of which 6 act before a score exists and 2 from inside a run. \textbf{Each is an existence
proof with its own denominator, and no incidence rate over models or platforms is computed from them.}

\textbf{The removals are not distributed evenly across the three platforms, and both denominators are
therefore reported.} Databricks 2 removals of 8 pairs attempted, Bedrock 3 of 10, \textbf{Azure 3 of 3}: the
overall 8 of 21 becomes 5 of 18 on the two measured platforms, 27.8\% against 38.1\%, a difference of 10.3
points resting on three attempts against one infrastructure. Three mechanisms below were observed only
there, which is why Table~\ref{tab:mechanisms} names the platform for each. Azure stays in the denominator
because dropping it would report a smaller number by discarding the attempts that failed most, and both
fractions are properties of an assembled roster rather than of managed clouds
(\S\ref{sec:proposizioni}).

Building the sample required probing which models each infrastructure would serve, and that probe returned
six mechanisms. Two more appeared from inside the confirmatory arm, where the roster probe does not reach
them: eight in total, none sought.

\subsection{The eight mechanisms, with the endpoint's own words}

\textbf{Every row is bounded to the configuration that produced it}, listed per mechanism in the deposit:
one endpoint, one API surface, one region, one set of parameters. The wording throughout is that the service
rejected the request under the evaluated configuration, not that the platform does not support the
feature. Where
variants were tried --- two identifier forms on M4, three history shapes on M6, the three admissible
tool-choice values on M7 --- the deposit records them, and those three rows can be stated more sharply for
that reason.

\begin{table}[!t]
\caption{The eight mechanisms, each an existence proof with its own denominator rather than a rate.
\textbf{Class} is operational and says what kind of thing the removal is, which is what decides who can do
anything about it: \emph{lifecycle} is an access or retirement policy, \emph{contract} is a constraint of
the API's own request contract, \emph{quota} is an inconsistency between what a catalogue advertises and
what the quota system holds, \emph{adapter} is an incompatibility between the format a model emits and the
validation the serving layer applies to it. \textbf{Scope} says how far the decision reaches:
\emph{vendor} holds for everyone on that platform, \emph{region} depends on where the request originates,
\emph{tenant} depends on the account's own history and quota, so two teams on one platform see different
rosters. \textbf{Det.} is the probe that reaches the mechanism, from the matrix in the deposit:
\emph{1} an ordinary single-request capability check, \emph{2} a deliberately constructed conversational
state, \emph{3} detectable only when the model itself emits the triggering output. Quoted cells are the
endpoint's words under the evaluated configuration, truncated to one line; M3 is a three-step control-plane
observation and is set unquoted for that reason. The deposit carries every message in full with the request
and the configuration that produced it.}
\label{tab:mechanisms}
\centering
\small
\setlength{\tabcolsep}{3pt}
\begin{tabular}{@{}l P{0.155\linewidth} l l l c P{0.375\linewidth}@{}}
\toprule
& mechanism & platform & class & scope & det. & verbatim message \\
\midrule
M1 & protocol refusal & Databricks & contract & vendor & 2 & \texttt{400}: multi-turn tool calls unsupported \\
M2 & deprecation & Azure & lifecycle & vendor & 1 & \texttt{ServiceModelDeprecated \ldots\ since 06/13/2026} \\
M3 & absent quota & Azure & quota & tenant & 1 & catalogue advertises it, quota system has no row, creation refused \\
M4 & disuse & Bedrock & lifecycle & tenant & 1 & \emph{``marked by provider as Legacy and you have not been actively using the model''} \\
M5 & jurisdiction & Bedrock & lifecycle & region & 1 & \emph{``Access to Meta Llama models is not allowed from unsupported countries''} \\
M6 & parallel calls & Azure & contract & vendor & 2 & \texttt{400 UnsupportedToolUse}: \emph{``does not support more than one tool call''} \\
M7 & tool config not expressible & Bedrock & contract & vendor & 2 & tool-choice admits ``auto'', ``any'', a named tool --- not ``none'' \\
M8 & endpoint rejects its own model's output & Bedrock & adapter & vendor & 3 & \texttt{ValidationException}: tool name fails \texttt{[a-zA-Z0-9\_-]+} \\
\bottomrule
\end{tabular}

\vspace{0.3em}
{\footnotesize Detection level 2 means a probe exists and must be built: for M1 a history already carrying
one tool call and its result; for M6 a history carrying two calls inside a single assistant message; for M7
a history carrying a tool-use block in a request that offers no tools. Each is one request once assembled.
Level 3 admits no such probe: the request must carry a tool name the model itself generated.}
\end{table}

Table~\ref{tab:mechanisms} is not a taxonomy derived from documentation: every row is a message an endpoint
returned to a request we made. M1--M6 were found by probing which models each infrastructure would serve;
M7 and M8 appeared from inside the confirmatory arm.

\textbf{The detection column is the operational finding, and it is a gradient rather than a split.} Four
mechanisms answer a single ordinary request. Three --- M1, M6 and M7 --- are invisible to any check that
sends one call per turn, and become one request each only once the conversational state is assembled by
hand; a harness certified healthy by an ordinary capability check can therefore still lose a cell after the
money is spent. M8 admits no probe at all in this sense, because what triggers it is generated by the model
during the run. The practical consequence for a pre-flight is that it must exercise realistic
conversational states, and that even then one class of failure remains reachable only by collecting.

For M6 the three-request sequence below separates the levels directly, and M7 is defined by a request
shape only a history containing a tool call can produce; M1's reproducer is the agentic exchange its
message names.

\subsection{Commentary on M6, M7 and M8}

\textbf{On M6.} The state it needs --- a history in which the model has already emitted two calls inside
one assistant message --- is one no single request constructs, but it \emph{is} constructible, and
assembling it by hand turns the mechanism into a probe answerable before any collection is paid for. Its
verification was three requests: one tool at turn one passes, a history with one prior tool call passes, a
history with two calls in one assistant message is refused.

\textbf{On M7.} The protocol we pre-registered ends each trajectory with a turn that offers no tools, and
under the evaluated Bedrock Converse configuration that turn \emph{cannot be expressed}: the API requires a
tool configuration whenever the history contains a tool call, and its tool-choice parameter admits
``auto'', ``any'' and a named tool --- but not ``none''. We tried the three admissible values and omitting
the configuration entirely; none expresses the turn.

\textbf{M7 is the easiest of the three level-2 mechanisms to have found, and we did not find it that
way.} The constraint is in the vendor's public API reference, which documents the tool-choice parameter as
a union of exactly three members with no value for ``no tools'': it was knowable before this study existed,
by reading, without issuing a request. M1 and M6 have no such public statement that we could locate, and
were reached by probing. The operative lesson is therefore not that M7 was undiscoverable but that
\emph{nobody checked the pre-registered protocol's termination shape against the target API's documented
request contract} --- a check that costs one page of reading and would have saved a 453-row cell.

The exposure was not where it first appeared. The broken rows carried an infrastructure-failure flag and
were already excluded from every mean; what survived was structural. A trajectory reaches the final turn
with a tool call in its history only if it has \emph{not} already submitted through a tool, so the runs
that completed were exactly those that had submitted early --- a second selection effect, inside the axis
this study exists to measure cleanly. That batch of 453 rows was discarded and the cell re-collected against
a corrected client. M7 sits at the same detection level as M6 and for the same reason: an ordinary
capability check reports it healthy, because it fires only in a conversational state that check does not
assemble.

\textbf{On M8.}
\label{sec:ottavo}
The open-weights model emits tool calls in its own chat format, which uses special tokens to separate
reasoning channels. The serving layer does not strip that token from the tool name, so the name reaches the
API as \texttt{decompile\_function<|channel|>commentary}. The API then validates the name it was just
handed:

\begin{quote}\footnotesize\ttfamily\raggedright
ValidationException: Value at\\
'messages.10.member.content.1.member.toolUse.name'\\
failed to satisfy constraint: Member must satisfy\\
regular expression pattern: [a-zA-Z0-9\_-]+
\end{quote}

\noindent
and rejects the entire request --- not the turn, the trajectory. \textbf{That sequence is what is reported here:
a name generated by the served model reaches the same service's request validation and is rejected there.}
It occurred in 11 trajectories, four in the original collection and seven in the re-collection, all in the
one cell where that model meets that endpoint under the native transport. Retrying does not reliably
recover the run: on one binary the cell logged 15 attempts for 7 valid measurements, one of the eight
rejected attempts carrying this exception verbatim and the other seven recording no candidate. Nothing in
the released evidence identifies where in the stack the format a model emits and the validation of the
service's own API fail to meet, and it is not assigned here.

Four further observations belong to the census and not to a paper of this length, and the deposit carries
each with its verbatim evidence. One carries an obligation that transfers: under a service control policy
that permits invocation while denying enumeration, the roster cannot be recovered from the account and must
be declared by explicit enumeration, which is what \S\ref{sec:design} does.

\subsection{How far a removal moves a published comparison}
\label{sec:bound}

That a refused model has no row to correct does not by itself say how much a reader's conclusion changes.
The census documents real cases --- a Llama variant deprecated on one platform while another serves it, a
second blocked by jurisdiction while its siblings run --- so someone evaluating on one endpoint reads a
table with rows missing and does not know which. We bound the consequence on the four models scored on both
clouds, removing one at a time and recomputing the comparison the way a leaderboard reports it. The
\emph{order} survives: all pairwise orderings agree before and after every removal. The \emph{scale} does
not, and with four models its collapse is close to arithmetic --- removing the lowest scorer necessarily
compresses a first-to-last spread --- so what these data support is the sign and the order of magnitude
rather than a coefficient. The artifact carries the eight removals with their per-cloud figures.

\subsection{What the census adds to what is already catalogued}
\label{sec:tassonomia}

The closest neighbour is a taxonomy of failure modes in multi-provider serving, classified by origin layer
and detectability~\citep{failureatlas}. What this census adds does not depend on where any single row is
filed in it, and that independence is deliberate: a reader who prefers to call M3 a configuration bug and
M8 an integration bug can do so, and four things stand unchanged.

\begin{enumerate}
\item \textbf{Multiple distinct removal mechanisms exist and were met in one small study} --- eight, none
  sought, each with the verbatim response of the endpoint that produced it.
\item \textbf{They sit at different points of the evaluation lifecycle.} Six act before any score exists,
  so the experimental unit is missing rather than mismeasured; two act inside a run that had already been
  paid for.
\item \textbf{Their measurement consequence is structural.} A removed pair has no row to correct, which is
  a selection and not a bias, and \S\ref{sec:bound} bounds what that does to a published comparison.
\item \textbf{They require different detection.} Four answer an ordinary capability check, three require a
  constructed conversational state, and one is reachable only during execution --- so a pre-flight that
  sends one request per model does not cover the class.
\end{enumerate}

\noindent
The question none of the existing layers asks is the one that decides whether a model appears in a
benchmark at all --- \emph{will this endpoint serve this model identity, under this request contract, from
this account and this region?} That is a reporting obligation before it is a taxonomic one, and it is where
Table~\ref{tab:checklist} puts it. There is no claim that nobody has catalogued serving failures; the artifact
carries the layer-by-layer comparison so a reader can place each row wherever they judge it belongs.
